\section{Thesis Objectives}
The studies performed in \cite{Cameron2013} and \cite{Schumann2015} provide a thought-provoking insight into the notion of aircraft plume superposition, and \cite{Dahlmann2020} demonstrates its potential utilisation for aviation climate impact mitigation through formation flight. This thesis aims to expand the scope of analysis conducted into chemical saturation effects that occur in high-density airspace. It intends to do so by satisfying the following objectives:
\begin{itemize}
    \item To conduct an expansive literature review into the state-of-the-art research surrounding aircraft emissions, aircraft exhaust plume dynamics, air traffic management and distribution, the local and global climate effects induced by aviation, and lastly, the mitigation of aviation climate impact, with a focus on operational strategies targeting non-\ce{CO_2} emissions.
    \item To utilise real-world air traffic data and first-order plume dispersion methods to assess the operational feasibility of aircraft exhaust plume overlap in today's airspace.
    \item To investigate the global spatial and temporal variability of \ce{NO_x} saturation effects, using a world-leading atmospheric chemistry transport model (CTM) STOCHEM-CRI, which incorporates the chemistry scheme CRIv2.2-R5 \cite{Khan2021}.
    \item To develop a modelling framework to robustly investigate the chemical and physical saturation effects that influence the non-\ce{CO_2} climate impact of formation flights.
    \item To carry out a localised study on aircraft \ce{NO_x} saturation effects using this modelling framework, in which parameters such as aircraft transit frequency and air temperature will be varied. 
\end{itemize}